\documentclass[11pt,a4paper]{article}
\usepackage[T1]{fontenc}
\usepackage[utf8]{inputenc}
\usepackage{lmodern}
\usepackage{amsmath,amssymb,amsthm,mathtools}
\usepackage[margin=28mm]{geometry}
\usepackage{microtype,booktabs,array,longtable}
\usepackage[hidelinks]{hyperref}
\hypersetup{pdftitle={Unbounded Boolean-Counting Obstructions in Endpoint Residue Packets}, pdfauthor={ChatGPT}}
\usepackage{fancyhdr}
\pagestyle{fancy}
\fancyhf{}
\fancyhead[L]{\small Endpoint residue obstructions}
\fancyhead[R]{\small ChatGPT}
\fancyfoot[C]{\thepage}
\setlength{\headheight}{14pt}
\setlength{\emergencystretch}{2em}
\newtheorem{theorem}{Theorem}[section]
\newtheorem{lemma}[theorem]{Lemma}
\newtheorem{proposition}[theorem]{Proposition}
\newtheorem{corollary}[theorem]{Corollary}
\theoremstyle{definition}
\newtheorem{definition}[theorem]{Definition}
\theoremstyle{remark}
\newtheorem{remark}[theorem]{Remark}
\DeclareMathOperator{\rad}{rad}
\DeclareMathOperator{\ord}{ord}
\newcommand{\FCRT}{B_{\mathrm{FCRT}}}
\newcommand{\SCRT}{B_{\mathrm{SCRT}}}
\newcommand{\Z}{\mathbb Z}
\newcommand{\N}{\mathbb N}
\newcommand{\pos}[1]{\left(#1\right)_{+}}
\title{Unbounded Boolean-Counting Obstructions\protect\\in Endpoint Residue Packets}
\author{ChatGPT}
\date{September 5, 2026}
\begin{document}
\maketitle
\begin{abstract}
We study a finite auxiliary inference in an endpoint-residue approach to the
standard $abc$ conjecture. A previous repository note exhibited a primitive
triple whose Boolean packet cube is larger than its residue group but has no
proper packet in the prescribed endpoint fibre. We extend that obstruction
to an infinite arithmetic family with an unbounded cube-to-group ratio.
For each $e\ge2$, Dirichlet's theorem and the Chinese remainder theorem give
squarefree products $B$ of $m=3^{e-1}$ distinct primes such that
$v_2(B+1)=2$, $v_3(B+1)=e$, and only the full packet is compatible with the
source $3$. A compatible saturated donor exists but has no such outgoing
proper-face flag. Nevertheless every constructed triple satisfies
$B+1<\rad(B(B+1))$ and has zero optimized flagged-CRT boundary. This
separates failure of an overly strong local rule from failure of a global
optimized estimate. We also give a standard antichain-based sufficient
criterion, a finite algebraic Lean core, a fixed-modulus arithmetic
formalization, and a deterministic exact replay. No proof or disproof of
standard $abc$, SCRT-0, or FCRT-1 is claimed.
\end{abstract}
\noindent\textbf{Keywords.} $abc$ conjecture; residue packets; zero-sum subsets;
Dirichlet's theorem; antichains; formal verification.

\section{Problem and scope}
For a positive integer $n$, let $\rad(n)$ be the product of its distinct prime
divisors. The standard target throughout the repository is
\begin{equation}\label{eq:abc}
 \forall\varepsilon>0\ \exists K_\varepsilon>0\ \forall a,b,c\in\N_{>0},\qquad
 \begin{gathered}a+b=c,\quad \gcd(a,b)=1\end{gathered}
 \Longrightarrow c\le K_\varepsilon\rad(abc)^{1+\varepsilon}.
\end{equation}
The quantifiers in \eqref{eq:abc} are not replaced by a bounded search, a
fixed-support theorem, or an alternative quality metric.

This note continues the repository at commit
\texttt{6118955}; the precise antecedents
are the endpoint-residue and anchored-prefix notes of September 4
\cite{repo}. Those notes already disprove the raw Boolean-counting shortcut
at $(1,4715,4716)$. Our contribution is an unbounded arithmetic extension
and its scope analysis, rather than a new refutation of an unexamined claim.
The construction gives no estimate for the difficult positive-defect region.
We make no claim of priority in the literature or of external peer review.

\subsection{Endpoint notation}
Fix a primitive positive triple $P=(a,b,c)$. In prime-indexed sets and
products, $p$ and $q$ always denote primes. Write
\[
 I=\{p\mid c:v_p(c)\ge2\},\qquad J=\{q:q\mid ab\},
 \qquad x_p=(v_p(c)-1)\log p,\quad y_q=\log q.
\]
Set $X=\sum_{p\in I}x_p$, $Y=\sum_{q\in J}y_q$. For $U\subseteq J$, take
whole prime-power factors on the two arms:
\[
 a_U=\prod_{\substack{q\in U\\q\mid a}}q^{v_q(a)},\qquad
 b_U=\prod_{\substack{q\in U\\q\mid b}}q^{v_q(b)}.
\]
If $p^{e_p}\Vert c$ with $e_p\ge2$, coprimality makes $a_U,b_U$ invertible
modulo $p^{e_p}$. Define
\begin{equation}\label{eq:label}
 G_p=(\Z/p^{e_p}\Z)^\times,\qquad
 \Phi_p(U)=a_U b_U^{-1}\in G_p.
\end{equation}
Thus $p^{e_p}\mid a_U+b_U$ if and only if $\Phi_p(U)=-1$. The full packet
always has this label. All subset products in this note are taken in the
abelian group in \eqref{eq:label}.

\subsection{The local flag and the global optimization}
For a nonempty source face $S\subseteq I$ and nonempty packet $T\subseteq J$,
put
\[
 M(S)=\prod_{p\in S}p^{e_p},\qquad x(S)=\sum_{p\in S}x_p,
 \qquad y(T)=\sum_{q\in T}\log q.
\]
The pair $(S,T)$ is a compatible block when $M(S)\mid a_T+b_T$, and it is
saturated when $x(S)\le y(T)$. It has surplus $\sigma=y(T)-x(S)$.
In the repository's FCRT-1 rule, a block can send one unsplit, one-hop token
to a residual source $p\notin S$ only if some nonempty proper $U\subsetneq T$
satisfies $\Phi_p(U)=-1$. Its credit is
\begin{equation}\label{eq:cap}
 \rho=\min\{\sigma,y(U)\}.
\end{equation}
Selected blocks have disjoint source faces and disjoint sink packets;
remaining sinks have exclusive residual owners. The optimized clipped
residual is denoted by $\FCRT(P)$. Discarding all surplus gives $\SCRT(P)$.
The previously proved accounting bound is
\begin{equation}\label{eq:bridge}
 \Delta(P):=\pos{\log c-\log\rad(abc)}\le\FCRT(P)\le\SCRT(P).
\end{equation}
The corresponding uniform upper estimate would imply \eqref{eq:abc}.
It is not established here. A prescribed donor without an outgoing flag
need not be part of an optimizing configuration.

\section{Elementary order calculation}
\begin{lemma}\label{lem:order}
For every integer $e\ge1$,
\[
 v_3\bigl(2^{3^{e-1}}+1\bigr)=e,
 \qquad \ord_{3^e}(2)=2\cdot3^{e-1}.
\]
\end{lemma}
\begin{proof}
Put $A_t=2^{3^t}$. Initially $A_0+1=3$. If $A_t=-1+3u$, then
\[
 A_t^2-A_t+1=3(1-3u+3u^2)
\]
has exact $3$-adic valuation one. The factorization
\[
 A_{t+1}+1=(A_t+1)(A_t^2-A_t+1)
\]
therefore proves $v_3(A_t+1)=t+1$ by induction.

For the order statement, take a positive exponent $n$ with
$2^n\equiv1\pmod{3^e}$. This first implies that $n$ is
even. Write $n=2u$ and $u=3^t s$ with $3\nmid s$. The factorization
\[
 4^{3^t}-1=(2^{3^t}-1)(2^{3^t}+1)
\]
has valuation $t+1$, since its first factor is not divisible by three.
The geometric sum
\[
 \frac{4^{3^t s}-1}{4^{3^t}-1}
   =\sum_{j=0}^{s-1}4^{3^t j}
\]
is congruent to $s$ modulo three, hence has valuation zero. Consequently
$v_3(4^u-1)=1+v_3(u)$. Divisibility by $3^e$ requires $3^{e-1}\mid u$.
The first part supplies the exponent $2\cdot3^{e-1}$, proving exactness.
\end{proof}

\section{An unbounded arithmetic obstruction}
\begin{theorem}\label{thm:main}
For each $e\ge2$, let $m=3^{e-1}$. There are infinitely many squarefree
integers $B$, each with exactly $m$ distinct prime factors, for which the
primitive triple $(1,B,B+1)$ has the following properties:
\begin{enumerate}
\item $v_2(B+1)=2$ and $v_3(B+1)=e$.
\item $|G_3|=2m$. Every sink label equals $2^{-1}$ in $G_3$; the full
packet has label $-1$, but no proper packet has label $-1$ and no nonempty
packet has label one.
\item The source face $\{2\}$ and the full sink packet form a compatible
saturated block of positive surplus, but no proper-face flag can target $3$.
\item The Boolean-cube/group-size ratio is $2^m/(2m)$ and is unbounded as
$e$ increases.
\end{enumerate}
\end{theorem}
\begin{proof}
Let $Q=3^{e+1}$. Since $8$ and $Q$ are coprime, the Chinese remainder theorem
gives a residue $r$ modulo $8Q$ such that
\begin{equation}\label{eq:crt}
 r\equiv2\pmod Q,\qquad r\equiv3\pmod8.
\end{equation}
This residue is coprime to $8Q$. Dirichlet's theorem gives infinitely many
primes in this progression; the same standard input is available in
Mathlib \cite{dirichlet}. Choose $m$ distinct such primes $q_1,\ldots,q_m$
and put $B=q_1\cdots q_m$. Varying the last prime above every bound proves
infinitude for each fixed $e$. The integer $B$ is squarefree and the triple
is primitive because its first coordinate is one.

The number $m$ is odd, so $B\equiv3^m\equiv3\pmod8$. Therefore
$v_2(B+1)=2$. Also
\[
 B+1\equiv2^m+1\pmod{3^{e+1}}.
\]
Lemma~\ref{lem:order} gives exact valuation $e$ for the right side and hence
for $B+1$. The modulus $3^{e+1}$, rather than $3^e$, is necessary for this
exact-valuation conclusion.

Every prime occurs to the first power on the $B$-arm. Its label is
$g=2^{-1}\pmod{3^e}$. By Lemma~\ref{lem:order}, $g$ has order $2m$ and
$g^m=-1$. A packet of cardinality $k$ has label $g^k$. If $0\le k<m$, this
cannot equal $g^m$ because exponents in $[0,2m)$ give distinct powers.
If $0<k\le m$, it cannot equal one either.

For the donor source $\{2\}$, the full modulus is four and the demand is
$\log2$. The full sink packet is compatible and has capacity $\log B>\log2$.
Its surplus is positive, but the proper target face necessary for a flag to
$3$ does not exist. Indeed, a flag from any block to $3$ would require a
proper subset of $J$ with target label $-1$ and is therefore impossible.

Finally $|(\Z/3^e\Z)^\times|=2\cdot3^{e-1}=2m$. For $m\ge3$,
\[
 \frac{2^m}{2m}\ge\frac{\binom m3}{2m}
 =\frac{(m-1)(m-2)}{12},
\]
which tends to infinity along $m=3^{e-1}$.
\end{proof}

\begin{corollary}\label{cor:constant}
For every real $C>0$, some primitive endpoints with a compatible saturated
donor satisfy $2^{|J|}>C|G_3|$ but have no proper target-compatible packet.
\end{corollary}
\begin{proof}
Choose $e$ making the ratio in Theorem~\ref{thm:main} exceed $C$.
\end{proof}

\begin{remark}
All equal-label fibres in the construction are rank layers: two subsets
have the same label precisely when they have the same cardinality.
There are very many collisions, but they are between incomparable sets.
Their signed differences are relations in the group; they are not positive
identity-product deletions. This is the exact obstruction to substituting a
raw Boolean pigeonhole argument for a nested-packet argument.
\end{remark}

\section{The optimized boundary is zero}
\begin{proposition}\label{prop:zero}
Every point constructed in Theorem~\ref{thm:main} satisfies
\[
 c<\rad(abc),\qquad \Delta(P)=0,\qquad
 \FCRT(P)=\SCRT(P)=0.
\]
\end{proposition}
\begin{proof}
Squarefreeness and $\gcd(B,B+1)=1$ give
\[
 \rad(abc)=B\rad(B+1)\ge6B>B+1=c,
\]
since both two and three divide $B+1$. This proves the first two claims.
For completeness, prime factorization gives the exact identities
\[
 X=\log\frac{c}{\rad(c)},\quad Y=\log B,\quad
 X-Y=\log c-\log\rad(abc)<0.
\]
The full block $(I,J)$ is compatible: its modulus divides $c=1+B$.
It is saturated because $X<Y$ and covers every source, leaving no residual.
The boundary is nonnegative by definition, so both optimized boundaries
are zero.
\end{proof}

The distinction between Theorem~\ref{thm:main} and
Proposition~\ref{prop:zero} is essential. The family refutes the assertion
that \emph{every prescribed saturated donor} has a useful outgoing flag.
It does not refute the existence of a good \emph{chosen configuration}.
In particular it supplies neither the fixed-positive-$\varepsilon$
counterfamily required to refute \eqref{eq:abc}, nor a counterfamily to the
optimized FCRT-1 uniform estimate.

The valid anchored-prefix criterion of the previous note is untouched.
Every nonempty reservoir in our construction generates a group of order
$2m$ but has at most $m$ generators. Its sufficient prefix-count hypothesis
therefore fails. The relevant research task is to find useful configurations
on the positive-defect region $X>Y$, not to insist on this stronger local
condition at arbitrary easy points.

\section{A valid positive antichain criterion}
The following standard combinatorial argument replaces the invalid raw-count
inference by a sufficient condition with a proved bound. It is not presented
as a new global estimate.

\begin{lemma}\label{lem:sperner}
An antichain of subsets of an $m$-element set has at most
$\binom{m}{\lfloor m/2\rfloor}$ members.
\end{lemma}
\begin{proof}
Each ordering of the ground set gives a maximal chain, for a total of $m!$
chains. An $r$-element subset belongs to exactly $r!(m-r)!$ such chains.
An antichain $\mathcal A$ meets any chain at most once, hence
\[
 \sum_{A\in\mathcal A}\frac{1}{\binom m{|A|}}\le1.
\]
Every binomial coefficient in the denominators is at most the central one,
which gives the asserted cardinality bound.
\end{proof}

\begin{theorem}\label{thm:antichain}
Attach labels in an abelian group to the elements of an $m$-element set $T$.
Suppose every subset product belongs to a finite subgroup of order $h$. If
\begin{equation}\label{eq:antichain}
 2^m>h\binom m{\lfloor m/2\rfloor},
\end{equation}
then a nonempty subset $D\subseteq T$ has product one. If the full packet
has label $\tau\ne1$, then $T\setminus D$ is a nonempty proper packet with
label $\tau$.
\end{theorem}
\begin{proof}
Assume no nonempty identity-product subset exists. A label fibre is an
antichain: if $A\subsetneq B$ have equal products, cancellation gives the
nonempty identity-product subset $B\setminus A$. Lemma~\ref{lem:sperner}
bounds each fibre by the central binomial coefficient. There are at most
$h$ fibres, contradicting \eqref{eq:antichain}. For the complement statement,
$D$ cannot be all of $T$ because the full label is not one; cancellation
then gives the target label for its nonempty proper complement.
\end{proof}

\begin{corollary}\label{cor:weighted}
Let a block $(S,T)$ be jointly compatible with a residual target $p\notin S$, with surplus
$\sigma\ge0$. Suppose a reservoir $K\subseteq T$ satisfies
\eqref{eq:antichain} with its own size and generated subgroup, and suppose
$y(K)\le x(S)$. Then an FCRT proper-face flag to $p$ reuses all of $\sigma$.
\end{corollary}
\begin{proof}
The theorem gives a nonempty deletion $D\subseteq K$ with identity label.
Its weight is at most $y(K)\le x(S)$. The complement $U=T\setminus D$
has target label and satisfies
\[
 y(U)=y(T)-y(D)\ge y(T)-x(S)=\sigma.
\]
The cap in \eqref{eq:cap} therefore equals $\sigma$.
\end{proof}

The extra width factor in \eqref{eq:antichain} prevents the raw counting
mistake. The criterion can be weaker than a useful ordered light-reservoir
argument, and no theorem here forces its weighted hypotheses on actual
positive-defect endpoints.

\section{Formalization and exact replay}
\subsection{Finite algebraic kernel}
The companion \texttt{EndpointConstantResidueObstruction.lean} contains
nine public theorems. It proves injectivity on the first half of a cyclic
orbit, the absence of identity-product nonempty packets, the absence of
proper packets in the full target fibre, and equal-rank collisions. The
hypothesis $\ord(g)=2m$ is explicit. The final declaration invokes Mathlib's
established Dirichlet theorem rather than adding an arithmetic axiom.
The relevant order API is documented in \cite{order}.

This core was strictly compiled under Lean 4.32.0 on the research commit
\texttt{c09e57f}. The dedicated PR run
\texttt{33956603573} completed successfully. The nine printed axiom reports
have union
\[
 \{\texttt{propext},\ \texttt{Classical.choice},\ \texttt{Quot.sound}\}.
\]
No global abc or uniform FCRT hypothesis occurs in these theorem types.

\subsection{A fixed-modulus arithmetic specialization}
The further companion \texttt{EndpointResidueArithmeticFamily.lean}
expresses an unbounded family without any abstract group-order premise.
Fix the primes $83$ and $947$, both congruent to $83$ modulo $216$, and take
an arbitrarily large prime
\[
 p\equiv83\pmod{216},\qquad p>947.
\]
Set $B=83\cdot947\cdot p=78601p$. Then $B\equiv35\pmod{216}$, giving
\[
 4\mid B+1,\quad8\nmid B+1,\quad9\mid B+1,\quad27\nmid B+1.
\]
Every factor is $2$ modulo nine. A proper packet has cardinality zero, one
or two, so its product is respectively $1$, $2$ or $4$ modulo nine; adding
one never gives zero. Dirichlet's theorem supplies $p$ above every bound.
These are precisely the arithmetic assertions of the specialization.
Its execution status is recorded in the accompanying verification report,
not inferred from the successful build of the earlier algebraic module.

The full $e$-uniform CRT/valuation realization, the logarithmic optimization
in Proposition~\ref{prop:zero}, and the antichain argument are ordinary
proofs in this manuscript. They are not labelled fully formalized merely
because one companion file compiles.

\subsection{Deterministic finite evidence}
The included Python program has no third-party dependencies. It selects the
first required primes in \eqref{eq:crt}, proving primality by trial division
through the integer square root. It checks exact $2$- and $3$-adic
valuations and the complete modular orbit. As the label depends only on
cardinality, checking all possible cardinalities covers every subset label
without enumerating all $2^m$ subsets.

\begin{center}
\begin{tabular}{rrrrr}
\toprule
$e$ & $m$ & $|G_3|$ & Largest chosen prime & Digits of $B+1$\\
\midrule
2 & 3 & 6 & 1,163 & 8\\
3 & 9 & 18 & 16,931 & 33\\
4 & 27 & 54 & 195,131 & 133\\
5 & 81 & 162 & 1,919,459 & 473\\
6 & 243 & 486 & 23,608,667 & 1,675\\
\bottomrule
\end{tabular}
\end{center}
The first triple is $(1,91412963,91412964)$. All five rows pass the exact
checks. The JSON digest is
\begin{center}\small\ttfamily
 ea56c8bc568e90311204e69e479534c9a7d\allowbreak
 afa2214db2c3c060831434fee9621.
\end{center}
The output's zero optimized boundary is derived from the full-block proof,
not from a numerical optimizer. The computation does not factor the large
number $B+1$ and does not test a uniform abc estimate.

\section{Relation to the other research routes}
The new result neither supplies nor invalidates the source construction
required by the repository's IUT program. Mochizuki's April 2026 report
\cite{iut} describes a five-stage formalization plan and its then-early
skeletal Stage 1. The report is not itself a released unconditional Lean
term for \eqref{eq:abc}. At the repository baseline, the conditional
\texttt{FourStageProgram} implication must still be distinguished from an
inhabitant built out of the stated arithmetic source objects \cite{repo}.

Likewise Bae's September 3 revision \cite{bae} gives an anchored-fibre
selection lemma in a binomial-coefficient problem. It can be applied with
its actual premises. Its application-specific tail estimates do not prove
an endpoint-reservoir condition, and our obstruction explains one failed
way of bypassing that missing step.

The remaining parallel gates in the repository are not closed by a finite
counting lemma. In a polynomial/Mason route, specialization and coefficient
heights must still be controlled. In the Frey/Szpiro, Vojta and S-unit routes,
constants must be uniform as the curves or supports vary. In Pell or
primitive-divisor routes, existence of a primitive divisor does not itself
supply exponent-one information. These are recorded as unresolved inputs,
not reasons to discard their parent routes.

\section{Conclusion}
A Boolean cube may overwhelm an endpoint group by an arbitrarily large
factor while containing no proper target-compatible packet. The phenomenon
occurs on genuine primitive arithmetic triples with a valid saturated donor,
not merely in an abstract group model. At the same time, every point in the
construction is harmless for the optimized global problem: the full block
has boundary zero. The next uniform theorem must therefore be about a
well-chosen configuration in the positive-defect region and must respect
nestedness and logarithmic weight, rather than infer a target face from raw
cardinality alone. No unconditional closed term of type
\texttt{ABCConjecture}, and no rigorous negation of that statement, results
from this work.

\begin{thebibliography}{9}
\raggedright
\bibitem{repo} \texttt{wyx66624/ABCConjecture}, research repository,
commit \texttt{6118955}, September 4, 2026.
In particular, \emph{Endpoint residue cubes and proper-subface flagged CRT
surplus}, \emph{Anchored prefix fibres and weighted proper-face CRT flags},
and \emph{ABC route bottlenecks}.\newline
\url{https://github.com/wyx66624/ABCConjecture}
\bibitem{dirichlet} The Mathlib Community, \emph{Primes in arithmetic
progression}, theorem \texttt{Nat.forall\_exists\_prime\_gt\_and\_modEq}.
\url{https://leanprover-community.github.io/mathlib4_docs/Mathlib/NumberTheory/LSeries/PrimesInAP.html}
\bibitem{order} The Mathlib Community, \emph{Order of an element}, theorem
\texttt{pow\_injOn\_Iio\_orderOf}.
\url{https://leanprover-community.github.io/mathlib4_docs/Mathlib/GroupTheory/OrderOfElement.html}
\bibitem{bae} J.~H.~Bae, \emph{Unbounded logarithmic limsup in Erd\H{o}s
Problem 684 via shifted carry scheduling}, arXiv:2604.23784v3,
September 3, 2026. \url{https://arxiv.org/html/2604.23784v3}
\bibitem{iut} S.~Mochizuki, \emph{Formalization of IUT}, workshop slides,
April 2026. \url{https://www.kurims.kyoto-u.ac.jp/~motizuki/Formalization%20of%20IUT%20(2026-04).pdf}
\end{thebibliography}
\end{document}
